\chapter*{TÓM TẮT KHOÁ LUẬN}
\addcontentsline{toc}{section}{{\bf TÓM TẮT KHOÁ LUẬN}\rm}

Khóa luận nghiên cứu và xây dựng một trợ lý ảo cho hệ thống e-learning dựa trên sinh tăng cường truy xuất (Retrieval-augmented generation - RAG). Trợ lý trả lời câu hỏi của người học từ học liệu đa định dạng gồm tài liệu PDF, DOCX, PPTX và video bài giảng, kèm trích dẫn có vị trí cụ thể (số trang hoặc thời điểm trong video) để người học kiểm chứng.
\bigbreak\noindent
Hệ thống thuộc hướng RAG đa phương thức (Multimodal RAG): học liệu được xử lý thành hai không gian vector, một cho văn bản và một cho ảnh, kèm một chỉ mục từ vựng BM25. Với video, hệ thống kết hợp nhận dạng giọng nói tự động (Automatic speech recognition - ASR), nhận dạng ký tự quang học (Optical character recognition - OCR) và chọn lọc khung hình đại diện. Ở bước truy xuất, hệ thống áp dụng truy xuất lai (hybrid retrieval) trên các nguồn văn bản, từ vựng và ảnh, hợp nhất bằng hợp nhất thứ hạng nghịch đảo (Reciprocal rank fusion - RRF) và lọc theo môn học. Câu trả lời được sinh kèm trích dẫn, và hệ thống từ chối trả lời khi bằng chứng không đủ tin cậy.
\bigbreak\noindent
Chúng tôi đánh giá khả năng truy xuất bằng một quy trình không rò rỉ dữ liệu trên bộ 400 câu hỏi về ba môn học, so sánh hệ thống đề xuất với một hệ thống cơ sở (baseline) chỉ truy xuất trên văn bản. Điểm được tính trên 325 câu hợp lệ, tương ứng độ phủ 85,53\% của kho học liệu đối với các câu có thể trả lời. Kết quả cho thấy chỉ khoảng 40\% số câu có ít nhất một kết quả đúng trong 10 kết quả đầu ở cả hai hệ thống, và hệ thống đề xuất chưa cải thiện nhất quán so với hệ thống cơ sở: khi lấy mẫu thống kê theo truy vấn, khác biệt không có ý nghĩa thống kê khi lấy 1, 5 hay 10 kết quả đầu, dù ở mức 5 kết quả đầu điểm của hệ thống đề xuất có xu hướng thấp hơn. Khóa luận chưa đánh giá chất lượng của câu trả lời được sinh; các thí nghiệm loại trừ (ablation study), đánh giá trên kho học liệu đầy đủ và đánh giá đầu-cuối (end-to-end evaluation) được ghi nhận là công việc tiếp theo.
